\chapter{\textbf{Formas Normales y la bifurcación de Takens-Bogdanov}}
\label{3:cap:formas_normales_tb}

El capítulo anterior introdujo SINDy como método data-driven cuya estructura matemática —regresión rala sobre una biblioteca de funciones— produce modelos parsimoniosos. Este capítulo da el paso siguiente: identifica el objeto teórico al que esa parsimonia apunta naturalmente —las \textbf{Formas Normales} cerca de bifurcaciones— y diseña el primer experimento de la tesis, que pone a prueba si SINDy efectivamente recupera ese objeto a partir de datos.

La elección del banco de prueba no es accidental. Las Formas Normales de bifurcaciones de codimensión 1 son demasiado simples como para distinguir un ajuste fenomenológico de la recuperación estructural: cualquier polinomio bajo orden las ajusta razonablemente. La codimensión 2, en cambio, presenta una geometría rica —múltiples curvas de bifurcación que se encuentran en un punto, regiones topológicamente distintas— que vuelve la prueba significativa. La bifurcación de \textbf{Takens-Bogdanov} (TB) es el caso prototípico: combina dos modos de inestabilidad en un solo punto degenerado y posee una forma normal polinómica con dos parámetros de control, lo que la hace particularmente compatible con el marco de SINDy presentado en el capítulo anterior.

El capítulo se organiza de la siguiente manera. La sección~\ref{3:sec:formas_normales} presenta el aparato de Formas Normales y la reducción a la variedad central. La sección~\ref{3:sec:codim_2} introduce el concepto de codimensión y motiva el estudio de codimensión 2. La sección~\ref{3:sec:tb} desarrolla la bifurcación de Takens-Bogdanov, su despliegue universal y la versión extendida usada como \textit{ground truth} en la tesis, junto con su diagrama de bifurcaciones y las cinco zonas dinámicas que lo dividen. Finalmente, la sección~\ref{3:sec:metodologia} presenta la metodología experimental: generación de datos, biblioteca de funciones, las cuatro estrategias de muestreo que estructuran los experimentos, y las métricas de evaluación con las que se diagnostica la recuperación.

\section{Formas Normales y reducción a la variedad central}
\label{3:sec:formas_normales}

\subsection{Idea y motivación}

Una Forma Normal es la representación más simple posible de un sistema dinámico cerca de un punto crítico, obtenida mediante cambios de coordenadas que eliminan todos los términos no esenciales del campo vectorial. La idea, debida originalmente a Poincaré y desarrollada de manera sistemática a lo largo del siglo XX, es que cerca de una singularidad la dinámica relevante está contenida en un puñado pequeño de monomios resonantes; el resto del campo vectorial se puede absorber en una transformación de coordenadas suficientemente regular sin cambiar el comportamiento cualitativo del sistema.

Formalmente, dado un sistema $\dot{x} = f(x)$ con un punto fijo en el origen, se buscan cambios de coordenadas locales $x = h(y)$, con $h$ difeomorfismo cercano a la identidad, tales que el sistema en las nuevas coordenadas $\dot{y} = g(y)$ retenga sólo los términos que no pueden ser eliminados por ningún $h$ de su clase de regularidad. Esos términos se llaman \textbf{resonantes}; el sistema $\dot{y} = g(y)$ es la Forma Normal asociada al punto fijo. La construcción explícita de las transformaciones $h$ procede orden por orden y está descripta con detalle, entre otras referencias, en los textos clásicos de Guckenheimer y Holmes~\cite{guckenheimer1983} y Kuznetsov~\cite{kuznetsov2004}.

La consecuencia práctica relevante para esta tesis es la siguiente: la Forma Normal es un objeto \textbf{ralo} por construcción. Cerca de una bifurcación, se reduce típicamente a un polinomio de bajo orden con unos pocos términos no nulos —los resonantes— mientras que el resto del campo vectorial original (típicamente de dimensión infinita en términos de su expansión en serie) queda absorbido en cambios de coordenadas. Esta rareza intrínseca es exactamente la estructura que el sesgo inductivo de SINDy busca, y constituye el argumento de fondo por el cual este trabajo apuesta a que el método sea adecuado para descubrir Formas Normales a partir de datos.

\subsection{Variedad central}

Cuando un sistema posee múltiples direcciones espectrales en un punto fijo —algunas estables, otras inestables, otras críticas—, el teorema de la Variedad Central permite reducir el análisis local a un subespacio de dimensión igual al número de autovalores con parte real nula. Las direcciones rápidas (con autovalores no nulos) se contraen o expanden exponencialmente y no participan de la dinámica esencial cerca de la bifurcación; toda la información cualitativa se concentra en la \textbf{variedad central}, una variedad invariante tangente al subespacio crítico en el punto fijo.

Este teorema, originalmente demostrado por Pliss y refinado por Kelley, Carr~\cite{carr1981} y otros, garantiza que para cualquier sistema $\dot{x} = f(x)$ con linealización admitiendo descomposición espectral $E^c \oplus E^s \oplus E^u$, existe localmente una variedad invariante $W^c$ tangente a $E^c$, sobre la cual la dinámica reducida captura toda la información topológica del sistema completo. La consecuencia, central para esta tesis, es que el estudio de bifurcaciones puede llevarse adelante en un sistema de baja dimensión —típicamente uno o dos— independientemente de la dimensionalidad ambiental del modelo original.

La combinación de variedad central y Formas Normales conforma así el aparato estándar de reducción local: primero se proyecta sobre la variedad central, eliminando los modos rápidos; luego se aplica el procedimiento de Formas Normales sobre el sistema reducido, eliminando los términos no resonantes. El resultado es una representación canónica de baja dimensión y baja cantidad de términos. Esta cadena de reducciones es la que vuelve compatible el problema con el marco de SINDy: tanto las dimensiones del estado como la cantidad de términos relevantes son chicas.

\subsection{Despliegues universales}

Las bifurcaciones, por definición, son fenómenos paramétricos: cambios cualitativos en la dinámica al variar uno o más parámetros de control. La extensión paramétrica de la teoría de Formas Normales se denomina \textbf{despliegue universal}: dada una singularidad —un punto fijo con cierta degeneración estructural—, su despliegue universal es la familia mínima de sistemas paramétricos que captura todos los modos en que la singularidad puede romperse genéricamente al perturbar el sistema.

El despliegue universal de una bifurcación de codimensión $k$ involucra exactamente $k$ parámetros independientes: si se perturba con menos parámetros, la perturbación es no genérica y no exhibe todos los regímenes posibles; si se perturba con más, la información adicional es redundante. Esta minimalidad paramétrica es la versión paramétrica de la rareza estructural de las Formas Normales, y refuerza el argumento de compatibilidad con SINDy: tanto el campo vectorial como su dependencia paramétrica son ralos en una base apropiada.

\section{Bifurcaciones de codimensión 2}
\label{3:sec:codim_2}

\subsection{Codimensión: definición operativa}

La \textbf{codimensión} de una bifurcación es el número mínimo de parámetros independientes que hay que variar para encontrarla genéricamente. De manera equivalente: es la cantidad de condiciones de no-degeneración que el campo vectorial debe satisfacer simultáneamente para que la bifurcación tenga lugar.

Las bifurcaciones de \textbf{codimensión 1} —saddle-node, transcrítica, pitchfork, Andronov-Hopf— ocurren genéricamente al variar un único parámetro; sus diagramas de bifurcación son curvas en el espacio de parámetros y sus Formas Normales son polinomios de pocos términos en una variable. Las de codimensión 2 requieren la variación simultánea de dos parámetros y se presentan como puntos aislados en el plano de parámetros, en los que confluyen varias curvas de bifurcación de codimensión 1.

La codimensión 2 es el primer escalón no trivial en la jerarquía de bifurcaciones. Su importancia es doble. Por un lado, cualquier sistema con dos parámetros de control —y la mayoría de los sistemas físicos lo son— exhibe genéricamente puntos de codimensión 2 en su diagrama; entender estos puntos es equivalente a entender la organización topológica global del diagrama de bifurcaciones. Por otro, la geometría es ya lo suficientemente rica como para que la recuperación de la Forma Normal a partir de datos sea un test informativo: distintas regiones del diagrama presentan distintos retratos de fase, y un método de descubrimiento que falle en distinguirlos quedará evidentemente expuesto.

\subsection{Catálogo breve}

Las bifurcaciones de codimensión 2 más estudiadas son:

\textbf{Cusp.} Doble degeneración en una bifurcación saddle-node: confluyen dos curvas de saddle-node en el plano $(\mu_1, \mu_2)$. La Forma Normal es $\dot{x} = \mu_1 + \mu_2 x \pm x^3$.

\textbf{Bautin (Hopf generalizada).} Una bifurcación de Hopf en la que el primer coeficiente de Lyapunov se anula. Curvas de Hopf super- y subcrítica se encuentran junto con una curva de bifurcación de ciclo límite saddle-node.

\textbf{Bogdanov-Takens.} Doble cero en la linealización: el Jacobiano tiene un autovalor doble en cero con bloque de Jordan no diagonalizable. De su despliegue emergen tres curvas de codimensión 1: saddle-node, Hopf y homoclínica.

\textbf{Hopf-Hopf, Fold-Hopf.} Dobles inestabilidades de tipo distinto, con dinámica considerablemente más compleja (toros invariantes, caos genérico).

Esta tesis se concentra exclusivamente en Takens-Bogdanov por dos razones: su despliegue universal es polinómico de bajo orden, lo que la vuelve perfectamente compatible con la biblioteca de SINDy; y su diagrama de bifurcaciones contiene cinco regiones topológicamente distintas, suficiente diversidad como para diseñar protocolos de muestreo no triviales.

\section{La bifurcación de Takens-Bogdanov}
\label{3:sec:tb}

\subsection{Definición y condiciones}

La bifurcación de Takens-Bogdanov ocurre en un punto fijo cuyo Jacobiano admite un autovalor doble en cero con bloque de Jordan no trivial. Concretamente, en una vecindad del punto crítico se puede llevar el sistema, por cambio de coordenadas lineal, a la forma

\begin{equation}
\frac{d}{dt}
\begin{pmatrix} x \\ y \end{pmatrix}
=
\begin{pmatrix} 0 & 1 \\ 0 & 0 \end{pmatrix}
\begin{pmatrix} x \\ y \end{pmatrix}
+
\begin{pmatrix} f(x,y) \\ g(x,y) \end{pmatrix},
\tag{3.1}
\label{eq:tb_lineal}
\end{equation}

donde $f, g$ son no lineales y se anulan al menos a segundo orden en el origen. Las condiciones de no-degeneración —que se desarrollan con detalle en el Capítulo 8 de Kuznetsov~\cite{kuznetsov2004}— exigen que ciertos coeficientes específicos de la expansión de Taylor de $f$ y $g$ sean no nulos; en particular, deben ser no nulos los coeficientes asociados a los monomios $x^2$ y $xy$ en el sistema reducido. Bajo estas condiciones, el despliegue universal canónico tiene la forma

\begin{equation}
\begin{aligned}
\dot{x} &= y, \\
\dot{y} &= \beta_1 + \beta_2 x + x^2 + s\, x y,
\end{aligned}
\tag{3.2}
\label{eq:tb_canonical}
\end{equation}

donde $s = \pm 1$ es un signo determinado por el sistema original y $(\beta_1, \beta_2)$ son los dos parámetros de despliegue. Para fijar ideas se trabaja con $s = -1$, que es la convención adoptada en la implementación de la tesis.

Del despliegue~\eqref{eq:tb_canonical} emergen, al recorrer el plano $(\beta_1, \beta_2)$, tres curvas de bifurcación de codimensión 1:

\begin{itemize}
    \item Una curva de \textbf{bifurcación saddle-node}, dada por $\beta_1 = \beta_2^2/4$, donde dos puntos fijos colisionan y se aniquilan.
    \item Una curva de \textbf{bifurcación de Hopf} —subcrítica para $s=-1$— que da nacimiento a un ciclo límite inestable.
    \item Una curva de \textbf{bifurcación homoclínica}, en la que el ciclo límite colisiona con un punto silla y se destruye.
\end{itemize}

Las tres curvas confluyen en el origen $(\beta_1, \beta_2) = (0, 0)$, que es el punto de Takens-Bogdanov propiamente dicho.

\subsection{Despliegue extendido y \textit{ground truth} de la tesis}

El despliegue canónico~\eqref{eq:tb_canonical} captura la dinámica local cerca del punto TB, pero no es globalmente bien comportado: el término $+x^2$ permite que las soluciones diverjan a infinito en tiempo finito si la condición inicial es suficientemente grande. Para los experimentos numéricos de esta tesis se requiere un sistema acotado globalmente, de modo que las trayectorias generadas como datos de entrenamiento no escapen del dominio durante la simulación. Por esa razón se trabaja con una versión \textbf{extendida} del despliegue, obtenida agregando términos estabilizadores de orden superior:

\begin{equation}
\boxed{
\begin{aligned}
\dot{x} &= y, \\
\dot{y} &= -\mu_1 - \mu_2\, x + x^2 - x^3 - x\, y - x^2\, y.
\end{aligned}
}
\tag{3.3}
\label{eq:tb_extendida}
\end{equation}

Esta es la ecuación implementada en el repositorio del trabajo y el \textit{ground truth} contra el cual se evalúan los modelos recuperados por SINDy. Los términos agregados respecto a~\eqref{eq:tb_canonical} son $-x^3$ —cúbico estabilizador en $\dot{y}$— y $-x^2 y$ —que junto al término $-xy$ vuelve la disipación dependiente del estado—; ambos son resonantes en el sentido de que aparecen genéricamente en el desarrollo extendido de la singularidad, y su presencia no altera la estructura topológica local del diagrama. La transformación de signos en los parámetros ($\beta_1 = -\mu_1$, $\beta_2 = -\mu_2$) es convencional. La tabla~\ref{tab:tb_coeffs} resume los coeficientes verdaderos del sistema, organizados por ecuación.

\begin{table}[h]
\centering
\begin{tabular}{c|c|c}
\hline
Ecuación & Término & Coeficiente verdadero \\
\hline
$\dot{x}$ & $y$       & $+1$ \\
\hline
$\dot{y}$ & $\mu_1$   & $-1$ \\
$\dot{y}$ & $x\mu_2$  & $-1$ \\
$\dot{y}$ & $x^2$     & $+1$ \\
$\dot{y}$ & $x^3$     & $-1$ \\
$\dot{y}$ & $xy$      & $-1$ \\
$\dot{y}$ & $x^2 y$   & $-1$ \\
\hline
\end{tabular}
\caption{Coeficientes verdaderos del despliegue extendido de Takens-Bogdanov utilizado en los experimentos. La ecuación tiene siete términos no nulos sobre una biblioteca polinómica de orden 3 en cuatro variables (ver §\ref{3:sec:metodologia}), que contiene 35 funciones candidatas: el problema es estructuralmente \textbf{ralo}.}
\label{tab:tb_coeffs}
\end{table}

\subsection{Diagrama de bifurcaciones y las cinco zonas dinámicas}

El plano de parámetros $(\mu_1, \mu_2)$ alrededor del origen queda dividido por las tres curvas de bifurcación —saddle-node, Hopf, homoclínica— en cinco regiones topológicamente distintas, que se reportan en la implementación del repositorio bajo el nombre genérico de \textbf{zonas}. La clasificación, basada en el número y la estabilidad de los puntos fijos y en la presencia de ciclos límite, es la siguiente:

\begin{description}
    \item[Zona 1.] Un único punto fijo, con $\mu_1 < 0$. El punto fijo es típicamente una silla.
    \item[Zona 2.] Un único punto fijo, con $\mu_1 \geq 0$. La estabilidad cambia respecto a la zona 1.
    \item[Zona 3.] Tres puntos fijos coexistentes, en la región $\mu_1 < 0$, $\mu_2 > 0$, separada de las zonas 1 y 2 por la curva saddle-node.
    \item[Zona 4.] Coexistencia de tres puntos fijos con un ciclo límite estable interno a la curva homoclínica, originado por la bifurcación de Hopf.
    \item[Zona 5.] Tres puntos fijos sin ciclo límite, posterior a la destrucción del ciclo en la bifurcación homoclínica (zona de escape).
\end{description}

La distinción entre Zona 4 y Zona 5 requiere localizar la curva homoclínica, que no admite forma cerrada: se determina numéricamente mediante un algoritmo de \textit{shooting} que detecta el valor de $(\mu_1, \mu_2)$ para el cual la variedad inestable de la silla se reconecta consigo misma~\cite{kuznetsov2004}. Las curvas saddle-node y Hopf, en cambio, son analíticamente conocidas y se calculan en forma cerrada a partir de la condición de doble raíz del polinomio $-x^3 + x^2 - \mu_2 x - \mu_1 = 0$ (saddle-node) y de la condición de traza nula del Jacobiano (Hopf).

La figura~\ref{fig:diagrama_tb} debe mostrar el plano $(\mu_1, \mu_2)$ con las tres curvas y las cinco zonas etiquetadas, junto a un retrato de fases representativo en cada zona. Esta figura es la referencia visual permanente del capítulo y del experimento subsiguiente.

\begin{figure}[h]
\centering
\fbox{\parbox{0.8\textwidth}{\centering\textit{[Insertar figura: diagrama de bifurcaciones de TB extendida en el plano $(\mu_1, \mu_2)$, con las curvas saddle-node (analítica), Hopf (analítica), homoclínica (detectada numéricamente) y las cinco zonas etiquetadas. Adjuntar paneles con retratos de fase representativos de cada zona.]}}}
\caption{Diagrama de bifurcaciones del despliegue extendido de Takens-Bogdanov.}
\label{fig:diagrama_tb}
\end{figure}

\section{Metodología: SINDy aplicado a Takens-Bogdanov}
\label{3:sec:metodologia}

\subsection{Generación de datos}

\textbf{Sistema de referencia.} Las trayectorias se generan integrando numéricamente la ecuación~\eqref{eq:tb_extendida} mediante un esquema de Runge-Kutta de orden 4 (RK4), implementado con \texttt{Numba} para compilación \textit{just-in-time} y ejecución a velocidad de C. La elección de RK4 frente a integradores adaptativos como Dormand-Prince busca un compromiso entre precisión y reproducibilidad: con paso fijo, el muestreo temporal es uniforme y la estimación posterior de $\dot{x}, \dot{y}$ por diferencias finitas no introduce sesgos derivados de la geometría adaptativa del paso.

\textbf{Grilla de parámetros.} Para cada experimento se define una región rectangular del plano $(\mu_1, \mu_2)$ y se la discretiza mediante producto cartesiano de \texttt{linspace} en cada eje (utilidad \texttt{generate\_param\_grid} en \texttt{core/utils.py}). La región estándar de exploración —denominada \texttt{base} en el código— es $\mu_1 \in [-0.1,\, 0.25]$, $\mu_2 \in [-0.1,\, 0.4]$, suficientemente amplia como para contener al punto TB en el origen y a las cinco zonas alrededor.

\textbf{Condiciones iniciales.} Para cada nodo de la grilla se integran $N_{\text{IC}}$ trayectorias con condiciones iniciales muestreadas de manera uniforme dentro de un rectángulo del espacio de fases acorde con la región paramétrica activa (atributo \texttt{state\_limits} del sistema). En los experimentos reportados, $N_{\text{IC}} = 10$ por nodo, lo que permite tanto promediar sobre regímenes transitorios como diagnosticar la sensibilidad a la condición inicial.

\textbf{Almacenamiento.} Las trayectorias resultantes se guardan en un único archivo HDF5 con compresión \texttt{gzip}, organizado por nodo de la grilla. Cada nodo se identifica con una clave de la forma \texttt{mu1\_X\_mu2\_Y}, y a cada clave se asocia un \textit{dataset} de forma $(N_{\text{IC}}, n, T)$ donde $n=2$ es la dimensión del estado y $T$ el número de pasos temporales. Junto a las trayectorias se guarda metadata por nodo (número de puntos fijos, estabilidades, métricas de calidad) que se usa más adelante para clasificar zonas y filtrar datos defectuosos.

\subsection{Biblioteca de funciones y aumento del estado}

Siguiendo la extensión paramétrica de SINDy descripta en el Capítulo~\ref{2:cap:identificacion_de_sistemas_no_lineales_mediante_representaciones_dispersas}, los parámetros $\mu_1, \mu_2$ se incorporan al estado como variables observables con derivada nula. La biblioteca de funciones candidatas es polinómica hasta orden 3 sobre las cuatro variables $(x, y, \mu_1, \mu_2)$, lo que da un total de

\[
\binom{4 + 3}{3} = 35
\]

funciones (constante, lineales, monomios cuadráticos y cúbicos, con todas las mezclas posibles entre estado y parámetros). El despliegue extendido~\eqref{eq:tb_extendida} contiene 7 términos no nulos sobre estos 35: la dispersión esperada del modelo recuperado es del orden del 80\,\%, lo que sitúa al problema bien adentro del régimen para el cual la regresión rala fue diseñada.

La elección del orden 3 se justifica por dos consideraciones. Primero, es el orden mínimo que contiene todos los términos del despliegue extendido (en particular $x^3$ y $x^2 y$). Segundo, mantiene el tamaño de la biblioteca manejable: subir a orden 4 sobre cuatro variables expande la biblioteca a 70 candidatos, lo que aumenta la probabilidad de identificar términos espurios sin agregar capacidad expresiva relevante para las Formas Normales de la familia TB.

\subsection{Estrategias de muestreo}
\label{3:subsec:casos_muestreo}

El experimento central del capítulo consiste en aplicar SINDy a cuatro escenarios distintos de selección de datos, diseñados para responder preguntas operativas progresivamente más realistas sobre lo que un experimentador puede esperar del método. Cada estrategia se resume en la tabla~\ref{tab:casos}.

\begin{table}[h]
\centering
\small
\begin{tabular}{c|p{4.2cm}|p{6.5cm}}
\hline
Caso & Datos disponibles & Pregunta que responde \\
\hline
\textbf{A} & Trayectorias en las cinco zonas, parámetros variando libremente. & ¿Recupera SINDy el despliegue extendido en el escenario más favorable posible? Cota superior del método. \\
\hline
\textbf{B} & Trayectorias confinadas a una sola zona; parámetros $(\mu_1,\mu_2)$ variando dentro de esa zona. & Laboratorio típico: el experimentador puede perturbar los parámetros de control pero no puede arrastrar al sistema a través de las curvas de bifurcación. \\
\hline
\textbf{C} & Trayectorias en una sola zona con parámetros fijos en un único valor de $(\mu_1, \mu_2)$. & Laboratorio sin control paramétrico: serie temporal larga con condiciones experimentales congeladas. \\
\hline
\textbf{D} & Trayectorias en una región del plano de parámetros lejana al punto TB (regiones \texttt{far\_z1}, \texttt{far\_z2}, \texttt{far\_z5} en el código). & Control negativo: ¿qué entrega SINDy si los datos no son cercanos a una degeneración? Caracterización del régimen \enquote{ajuste fenomenológico} por contraste con la recuperación estructural de los casos A--C. \\
\hline
\end{tabular}
\caption{Las cuatro estrategias de muestreo. Cada caso responde una pregunta operativa específica sobre el alcance del método.}
\label{tab:casos}
\end{table}

\textbf{Caso A.} Se entrena SINDy con un conjunto de trayectorias que cubre las cinco zonas. Operativamente, se selecciona aleatoriamente un \textit{batch} de claves de la grilla cuya unión cubra los cinco rótulos de zona, y se concatena el conjunto correspondiente de trayectorias —cada una etiquetada por su par $(\mu_1, \mu_2)$— en la matriz de datos aumentada. Es el escenario más favorable: la diversidad de regímenes topológicos en los datos es máxima.

\textbf{Caso B.} Se restringe el muestreo a una sola zona. Se repite el experimento independientemente para cada una de las cinco zonas. Esto permite diagnosticar si SINDy puede recuperar la Forma Normal sin ver todos los regímenes topológicos, apoyándose únicamente en la dependencia paramétrica local dentro de la zona observada. La hipótesis es que la información paramétrica es tan o más valiosa que la diversidad topológica para la identificación.

\textbf{Caso C.} Versión más restrictiva del Caso B: $(\mu_1, \mu_2)$ fijos en un único nodo de la grilla. Las trayectorias de entrenamiento corresponden todas al mismo punto en parámetros, con condiciones iniciales variables. La predicción es que SINDy pierde los términos del campo vectorial que dependen explícitamente de $\mu$ —ya que estos no varían en los datos— pero conserva la estructura no lineal en $(x, y)$ con coeficientes congelados en el valor que asumen para ese par específico.

\textbf{Caso D.} Se cambian los rangos paramétricos a una de las regiones \texttt{far} predefinidas, lejanas al punto TB. Las trayectorias generadas no exhiben la dinámica característica de la singularidad y SINDy debe recurrir a un ajuste polinómico local. La predicción es que el modelo recuperado no será estructuralmente reconocible como el despliegue extendido sino un ajuste válido localmente.

El conjunto de los cuatro casos produce el \textit{mapa de operación} del método: una caracterización empírica de en qué condiciones de muestreo se recupera la Forma Normal y en cuáles se obtiene un ajuste sin estructura canónica.

\subsection{Optimizador y selección de hiperparámetros}

Sobre la biblioteca polinómica se entrena el modelo mediante \textit{Sequentially Thresholded Least Squares} (STLSQ), implementado en la librería \texttt{PySINDy}~\cite{desilva2020pysindy}. Los hiperparámetros del optimizador son dos:

\begin{itemize}
    \item El \textbf{umbral de dispersión} $\lambda$ (\texttt{threshold}), que controla la corteza de la regresión rala: en cada iteración se anulan los coeficientes con magnitud menor a $\lambda$.
    \item El \textbf{coeficiente Ridge} $\alpha$ (\texttt{alpha}), que añade una penalización $\ell_2$ al problema de mínimos cuadrados que se resuelve sobre el soporte vigente, estabilizando la solución frente a multicolinealidad de columnas en $\Theta$.
\end{itemize}

En todos los experimentos reportados, las derivadas $\dot{x}, \dot{y}$ se estiman mediante diferencias finitas (\texttt{ps.FiniteDifference}) sobre las trayectorias muestreadas, y el modo \texttt{unbias} se evalúa en sus dos configuraciones (\texttt{True} y \texttt{False}). En la primera, los coeficientes finales se obtienen mediante un mínimos cuadrados sin regularización sobre el soporte sobreviviente —eliminando el sesgo hacia cero característico del LASSO—; en la segunda, se conservan los coeficientes obtenidos directamente del proceso iterativo de poda. La comparación entre ambos modos es relevante porque caracteriza la sensibilidad del método al post-procesamiento.

La búsqueda de hiperparámetros se hace por dos vías complementarias. Primero, una \textit{grid search} sobre $(\lambda, \alpha)$ retiene los $K$ mejores modelos según una métrica de error sobre los coeficientes (descripta en la subsección siguiente); este barrido se implementa en el script \texttt{run\_optimization.py}. Segundo, los modelos seleccionados se refinan localmente mediante \textit{hill climbing} sobre dos ejes: la composición del \textit{batch} de datos de entrenamiento (intercambiando claves al azar) y los hiperparámetros (perturbaciones multiplicativas del $\pm$\,20\,\% sobre $\lambda$ y $\alpha$). Esta segunda fase se implementa en \texttt{run\_fine\_tuning.py}.

\subsection{Métricas de evaluación}

La evaluación de un modelo recuperado se hace contra el \textit{ground truth} de la tabla~\ref{tab:tb_coeffs}, comparando el soporte y los valores numéricos de los coeficientes. Se emplean tres categorías de métrica, complementarias entre sí:

\textbf{Error relativo por término esperado.} Para cada término no nulo del despliegue verdadero $\xi^{\text{true}}_{ij}$, se calcula el error relativo
\[
e_{ij} = \frac{|\hat{\xi}_{ij} - \xi^{\text{true}}_{ij}|}{|\xi^{\text{true}}_{ij}|},
\]
donde $\hat{\xi}_{ij}$ es el coeficiente recuperado por SINDy. Un término \enquote{verdadero} que SINDy haya descartado por umbralado tiene $\hat{\xi}_{ij} = 0$ y contribuye con $e_{ij} = 1$ (etiqueta \texttt{MISSING} en el código).

\textbf{Términos espurios (\texttt{GHOST}).} Cualquier término no nulo en $\hat{\Xi}$ que no aparezca en el despliegue verdadero se reporta como espurio. Su presencia indica que el método identificó estructura inexistente, ya sea por sobreajuste a ruido, por correlación entre columnas de $\Theta$, o por insuficiente diversidad en los datos. Se cuantifica su número y la magnitud de sus coeficientes.

\textbf{Error promedio agregado.} La métrica única que se utiliza para rankear modelos en la grid search es el promedio sobre los términos esperados de $e_{ij}$, sumado a la magnitud de los términos espurios:
\[
E(\hat{\Xi}) = \frac{1}{N_{\text{true}}}\Bigg(\sum_{(i,j) \in \mathcal{S}_{\text{true}}} e_{ij} + \sum_{(i,j) \notin \mathcal{S}_{\text{true}}} |\hat{\xi}_{ij}|\Bigg).
\]
Esta combinación penaliza simétricamente los errores en términos verdaderos y los términos espurios, lo que evita que un modelo \enquote{trampee} la métrica anulando muchos coeficientes correctos a cambio de no introducir ningún espurio.

\textbf{Recuperación topológica.} Adicionalmente, se simula la dinámica del modelo recuperado y se compara cualitativamente con la del sistema verdadero: número de puntos fijos, estabilidades, presencia o ausencia de ciclos límite. Esta verificación es una salvaguarda final: pequeños errores numéricos en los coeficientes pueden o no preservar la estructura cualitativa del retrato de fases, y la comparación topológica permite diagnosticarlo.

\subsection{Barrido temporal estadístico}
\label{3:subsec:barrido_temporal}

Una decisión metodológica adicional que merece atención propia es el \textbf{barrido sobre la ventana de observación}. Para una configuración dada (zona, hiperparámetros, modo \texttt{unbias}), se entrena SINDy repetidamente truncando las trayectorias a $t \in [0, t_{\max}]$ con $t_{\max}$ recorriendo una grilla $\{t_1, \ldots, t_M\}$. Para cada $t_{\max}$ se ejecutan $N_{\text{boot}}$ corridas con muestreos bootstrap distintos del conjunto de trayectorias, y se reporta la media y desviación estándar de cada coeficiente recuperado a lo largo del barrido. Esto produce, por término candidato, una curva $(t_{\max}, \mu_{\hat\xi}(t_{\max}) \pm \sigma_{\hat\xi}(t_{\max}))$ que diagnostica la convergencia del método: idealmente, los términos verdaderos convergen a sus valores teóricos con varianza decreciente, mientras que los términos espurios se mantienen alrededor de cero o presentan varianza sistemáticamente alta.

En los experimentos reportados, $t_{\max}$ recorre una grilla densa entre 0.05\,s y 20.1\,s, y $N_{\text{boot}} = 10$. Esta elección permite observar la transición desde el régimen de dato escaso —en el que el sistema lineal está mal condicionado— hasta el régimen asintótico donde los coeficientes deberían estabilizarse en sus valores verdaderos.

El barrido temporal estadístico es la métrica principal de este capítulo. Su lectura permite distinguir tres situaciones cualitativamente distintas que se reportan en el capítulo siguiente:

\begin{enumerate}
    \item \textbf{Recuperación estable:} todos los términos verdaderos convergen a sus valores teóricos con varianza tendiente a cero; los espurios se anulan o quedan de varianza alta.
    \item \textbf{Recuperación parcial:} algunos términos verdaderos convergen, otros oscilan o convergen a valores incorrectos; aparecen espurios persistentes.
    \item \textbf{No recuperación:} el soporte recuperado no se corresponde con el despliegue verdadero, ni siquiera parcialmente.
\end{enumerate}

El objetivo de los experimentos es localizar, sobre el espacio de configuraciones (zona $\times$ caso de muestreo $\times$ hiperparámetros), las regiones en las que se obtiene cada uno de los tres regímenes. Los resultados se presentan en el Capítulo~\ref{X:cap:resultados_tb}.

\section{Cierre del capítulo}

Este capítulo estableció el marco teórico —Formas Normales, codimensión 2, bifurcación de Takens-Bogdanov— y la metodología experimental que sostiene el primer experimento de la tesis. Tres ideas conviene retener antes de pasar a los resultados.

Primero, la Forma Normal extendida~\eqref{eq:tb_extendida} es \textbf{rala por construcción}: 7 términos no nulos sobre una biblioteca de 35 candidatos. La compatibilidad estructural con SINDy es máxima.

Segundo, el diagrama de bifurcaciones presenta cinco zonas topológicamente distintas, lo que permite diseñar el espacio de experimentos como un cruce entre estrategia de muestreo (cuatro casos) y zona del diagrama (cinco rótulos), con un total no trivial de configuraciones a barrer.

Tercero, la metodología de evaluación combina métricas estructurales (soporte recuperado), numéricas (error en coeficientes) y temporales (convergencia bajo barrido de $t_{\max}$), de modo que ningún resultado del próximo capítulo puede interpretarse simplemente como \enquote{funcionó} o \enquote{no funcionó}: cada configuración produce un perfil de recuperación que admite lectura cualitativa.

Con este aparato montado, el Capítulo~\ref{X:cap:resultados_tb} reporta los resultados experimentales para cada uno de los cuatro casos de muestreo y discute en qué medida la Forma Normal de Takens-Bogdanov puede recuperarse a partir de datos.

\clearpage
